\uuid{NEvK}
\exo7id{3559}
\titre{exo7 3559}
\auteur{quercia}
\organisation{exo7}
\datecreate{2010-03-10}
\isIndication{false}
\isCorrection{false}
\chapitre{Réduction d'endomorphisme, polynôme annulateur}
\sousChapitre{Polynôme annulateur}
\cpge{mpsi-mp2i.calcul-matriciel-systemes.e, pcsi.determinants.c, ptsi.determinants.a, mpsi-mp2i.changements-bases.c, bcpst1.matrices, mp-mpi.reduction-endomorphismes.b, pc.reduction.a, psi.reduction-endomorphismes.a, pt.reduction-endomorphismes.a, mp-mpi.reduction-endomorphismes.d, pc.reduction.c, psi.reduction-endomorphismes.c, pt.reduction-endomorphismes.c, bcpst2.valeurs-propres-vecteurs-propres.a, bcpst2.valeurs-propres-vecteurs-propres.b}
\module{Algèbre}
\niveau{L2}
\difficulte{}

\contenu{
\texte{
Soit $M \in \mathcal{M}_{3}(\R)$ telle que $M \ne I$, et $M^3 = I$.
}
\begin{enumerate}
    \item \question{Quelles sont les valeurs propres complexes de $M$ ?
    (On vérifiera que ce sont effectivement des valeurs propres de $M$)}
    \item \question{Montrer que $M$ est semblable à
    $\begin{pmatrix} 1 &0    &0         \cr 
               0 &-1/2 &-\sqrt3/2 \cr
               0 &\sqrt3/2 &-1/2  \cr \end{pmatrix}$.}
\end{enumerate}
}
